\section{CONTEXTO Y MARCO TEÓRICO}

Este capítulo desarrolla los fundamentos conceptuales y técnicos que sostienen el análisis molecular y genómico de virus patógenos en aguas servidas. El marco teórico articula la vigilancia epidemiológica, la vigilancia basada en aguas residuales, los virus seleccionados, las técnicas de detección molecular, las plataformas de secuenciación y las herramientas bioinformáticas necesarias para interpretar señales virales procedentes de una matriz ambiental compleja \cite{Sims_2020,Parkins_2023,Child_2023}.

\subsection{Vigilancia epidemiológica}

\subsubsection{Importancia y aplicaciones}

La vigilancia epidemiológica permite transformar observaciones dispersas sobre eventos de salud en información útil para orientar decisiones sanitarias. En enfermedades infecciosas, su valor se expresa en la detección temprana de circulación, el seguimiento de tendencias, la evaluación de intervenciones y la priorización de recursos en poblaciones concretas. La vigilancia clínica aporta confirmación diagnóstica individual, pero depende de la búsqueda de atención, la disponibilidad de pruebas, la oportunidad de notificación y la capacidad del sistema de salud para sostener el monitoreo en el tiempo \cite{Parkins_2023,Bubba_2023}.

La vigilancia ambiental amplía este panorama al incorporar señales poblacionales obtenidas fuera del circuito clínico. Las aguas residuales reúnen material biológico excretado por personas sintomáticas, presintomáticas y asintomáticas conectadas a una red de alcantarillado, por lo que pueden aportar una lectura agregada de circulación microbiana en una comunidad. Esta aproximación resulta pertinente para patógenos con excreción fecal directa o indirecta, virus respiratorios detectables en aguas servidas y agentes con transmisión silenciosa o subregistro clínico \cite{Sims_2020,Parkins_2023}.

El seguimiento de SARS-CoV-2 durante la pandemia consolidó el uso de aguas residuales como herramienta complementaria para salud pública. La detección temprana de ARN viral en aguas residuales de Australia y Países Bajos mostró que la señal ambiental podía acompañar la circulación comunitaria cuando la vigilancia clínica aún estaba limitada por capacidad diagnóstica, acceso a pruebas y retrasos de reporte \cite{Ahmed_2020,Medema_2020}. Estudios posteriores demostraron correlaciones entre la carga viral en aguas residuales y la carga clínica de COVID-19 a escalas urbanas y submunicipales \cite{Acosta_2022,Sangsanont_2022}.

\subsubsection{Vigilancia basada en aguas residuales}

La vigilancia basada en aguas residuales consiste en el análisis sistemático de muestras de alcantarillado o influentes de plantas de tratamiento para identificar marcadores biológicos y químicos asociados con una población. Esta estrategia permite trabajar a diferentes escalas: plantas de tratamiento, estaciones de bombeo, redes barriales, instituciones o sitios centinela. La selección del punto de muestreo define la población representada, el nivel de resolución espacial y el tipo de decisiones que pueden apoyarse con los datos obtenidos \cite{Parkins_2023}.

En virus humanos, la vigilancia basada en aguas residuales ha tenido aplicaciones consolidadas en poliovirus y una expansión acelerada durante COVID-19. La vigilancia ambiental de poliovirus se desarrolló como complemento de la vigilancia de parálisis flácida aguda dentro de los esfuerzos de erradicación, debido a que puede detectar circulación viral en ausencia de casos paralíticos reportados \cite{Asghar_2014,Bubba_2023}. Para SARS-CoV-2, los estudios longitudinales en plantas de tratamiento y sectores urbanos demostraron que las concentraciones de ARN viral pueden reflejar cambios en transmisión comunitaria y anticipar aumentos observados por vigilancia clínica \cite{Acosta_2022,Sangsanont_2022}. En RSV, el ARN detectado en sólidos sedimentados se ha asociado con tasas de positividad clínica, lo que respalda su uso como señal ambiental de actividad respiratoria comunitaria \cite{Hughes_2022}.

Su interpretación exige considerar variables propias de la matriz y de la red sanitaria. El caudal, la dilución por lluvias o infiltración, la contribución industrial, el tamaño poblacional conectado, los tiempos de residencia, la estabilidad del ARN y la presencia de inhibidores pueden modificar la señal medida en laboratorio \cite{Parkins_2023}. La concentración previa de partículas virales, el uso de controles de extracción y la normalización con marcadores fecales o de proceso son medidas importantes para reducir incertidumbre analítica \cite{Hughes_2022,Child_2023}. Por esta razón, los resultados de aguas residuales deben leerse como indicadores poblacionales y metodológicamente condicionados, con mayor fortaleza para detectar presencia, tendencias y cambios relativos que para estimar casos clínicos individuales \cite{Parkins_2023,Bubba_2023}.

\subsubsection{Muestreo pasivo}

El muestreo pasivo se basa en la exposición de materiales adsorbentes o retentivos dentro del flujo de aguas residuales durante un periodo definido. Estos materiales capturan partículas virales, ácidos nucleicos o sólidos asociados, generando una muestra integrada en el tiempo sin requerir equipos automáticos de muestreo compuesto. Las revisiones sobre vigilancia en aguas residuales describen esta estrategia como una alternativa operativa para sitios con limitaciones de infraestructura, electricidad, seguridad o presupuesto \cite{Parkins_2023}.

En Quito, el antecedente local más directo es el trabajo de Mejia Calle en la PTAR Quitumbe. Ese estudio empleó un dispositivo pasivo tipo torpedo impreso en 3D, con membrana de nailon, gasa y hisopos de algodón, colocado durante 24 horas en el influente de la planta. El muestreo fue semanal, cubrió 47 muestras y representó aguas residuales procedentes de once barrios del Distrito Metropolitano de Quito, con un caudal reportado de 75 L/s \cite{MejiaCalle_2024}. Esta experiencia demuestra factibilidad local para recolectar material viral en un sitio centinela urbano y aporta una base metodológica para ampliar el análisis hacia varios patógenos.

La señal obtenida por muestreo pasivo depende del material utilizado, la duración de exposición, la velocidad del flujo, la carga de sólidos, la degradación del ARN, la afinidad del virus por la matriz y la eficiencia de recuperación en laboratorio. En el estudio de Quitumbe, la baja concentración viral en componentes individuales llevó a analizar componentes combinados del dispositivo, y la secuenciación fue condicionada por valores de Ct tardíos y por cambios en el rendimiento de cebadores \cite{MejiaCalle_2024}. Estas consideraciones son centrales para interpretar resultados positivos, negativos o de baja cobertura en un diseño multipatógeno.

\subsection{Virus de interés}

\subsubsection{SARS-CoV-2}

SARS-CoV-2 es un virus respiratorio de ARN cuya vigilancia ambiental adquirió relevancia desde las primeras etapas de la pandemia de COVID-19. La detección confirmada de ARN viral en aguas residuales sin tratar en Australia y Países Bajos respaldó la hipótesis de que el virus podía ser monitoreado mediante alcantarillado, debido a la eliminación de material genético viral por personas infectadas \cite{Ahmed_2020,Medema_2020}. Esta evidencia inicial fue seguida por estudios longitudinales que relacionaron la señal ambiental con indicadores clínicos y poblacionales de transmisión \cite{Acosta_2022,Sangsanont_2022}.

El valor de SARS-CoV-2 como blanco de vigilancia molecular se debe a tres características: alta relevancia sanitaria, disponibilidad de ensayos RT-qPCR validados y capacidad de caracterización genómica de variantes. Los genes N1 y N2 han sido utilizados de forma amplia como blancos de detección en aguas residuales, mientras la secuenciación permite identificar linajes y mezclas de variantes presentes en la población conectada a la red \cite{Ahmed_2020,Crits_Christoph_2021,MejiaCalle_2024}. En Quito, la PTAR Quitumbe permitió detectar señales de SARS-CoV-2 y asignar linajes de Ómicron mediante secuenciación Oxford Nanopore y análisis bioinformático con Freyja \cite{MejiaCalle_2024}.

La vigilancia genómica de SARS-CoV-2 en aguas residuales trabaja con una mezcla comunitaria de genomas virales. Esa mezcla puede revelar linajes predominantes, variantes emergentes y señales de baja frecuencia, aunque su interpretación depende de cobertura, profundidad, degradación del ARN, sesgos de amplificación y selección de referencias \cite{Crits_Christoph_2021,Child_2023,Parkins_2023}. En el antecedente de Quito, la detección de JN.1 en aguas residuales antes de su identificación clínica local muestra la utilidad de la matriz ambiental para observar cambios de linajes en contextos con vigilancia clínica genómica limitada \cite{MejiaCalle_2024}.

\subsubsection{Virus sincitial respiratorio (RSV)}

El virus sincitial respiratorio es un patógeno respiratorio asociado con carga importante de enfermedad en lactantes, adultos mayores e individuos inmunocomprometidos. Su vigilancia clínica puede verse limitada por subdiagnóstico, disponibilidad de pruebas, presentación clínica compartida con otros virus respiratorios y variaciones estacionales \cite{Hughes_2022,Le_2025}. La reciente disponibilidad de vacunas y estrategias de inmunización pasiva incrementa la necesidad de herramientas que permitan observar cambios en circulación y diversidad genética del virus \cite{Le_2025}.

La detección de RSV en aguas residuales ha sido documentada en sólidos sedimentados y en flujos de vigilancia genómica. Hughes y colaboradores observaron asociación entre ARN de RSV en aguas residuales y positividad clínica, con valores de Kendall tau entre 0,65 y 0,77 y significancia estadística alta \cite{Hughes_2022}. Le y colaboradores desarrollaron un flujo de amplificación y secuenciación de genoma completo con Oxford Nanopore, aplicable a muestras clínicas y aguas residuales, y reportaron que las frecuencias de clados observadas en aguas residuales representaban las tendencias observadas en pacientes \cite{Le_2025}.

Desde el punto de vista genómico, RSV presenta dos subtipos principales, RSV-A y RSV-B, con diversidad de clados que requieren referencias adecuadas para evitar pérdida de lecturas durante el mapeo. El flujo RSVTyper automatiza la identificación de subtipo, la selección de referencia y el análisis posterior, lo que mejora la interpretación de datos de secuenciación en muestras con mezclas o cargas variables \cite{Le_2025}. Para este estudio, RSV es relevante porque permite evaluar la factibilidad de incorporar un virus respiratorio de alta carga clínica dentro de un esquema ambiental originalmente aplicado de forma local a SARS-CoV-2.

\subsubsection{Poliovirus}

El poliovirus mantiene relevancia internacional por su relación con la erradicación global, el riesgo de reintroducción y la emergencia de poliovirus derivados de vacuna. La vigilancia clínica se centra en parálisis flácida aguda, un desenlace infrecuente respecto al total de infecciones. La literatura sobre vigilancia ambiental señala que la excreción viral puede detectarse en aguas residuales antes de observar casos paralíticos, lo que convierte al alcantarillado en una herramienta sensible para identificar circulación silenciosa \cite{Asghar_2014,Zurbriggen_2008,Bubba_2023}.

Los estudios ambientales han demostrado que la detección de poliovirus en aguas residuales puede ocurrir incluso en países con alta cobertura sanitaria. En Suiza, Zurbriggen y colaboradores aislaron poliovirus Sabin-like en muestras de aguas residuales y destacaron la necesidad de herramientas sensibles debido a la baja tasa de ataque clínico \cite{Zurbriggen_2008}. En 2024, la detección de cVDPV2 en aguas residuales de Finlandia, Alemania, Polonia, España y Reino Unido reveló cepas genéticamente relacionadas y activó una alerta sobre circulación o reintroducción en territorios con alta capacidad de vigilancia \cite{Bottcher_2025}.

La caracterización molecular es indispensable para diferenciar poliovirus Sabin-like, poliovirus salvaje y poliovirus derivados de vacuna. El análisis de la región VP1, la detección intratípica y la secuenciación permiten establecer relaciones genéticas, inferir origen y orientar respuestas de salud pública \cite{Zurbriggen_2008,Asghar_2014}. Shaw y colaboradores describen que los flujos tradicionales para confirmar brotes de VDPV2 pueden tomar semanas, mientras que la detección molecular directa y la secuenciación Nanopore tienen potencial para acortar tiempos de confirmación y respuesta \cite{Shaw_2022}.

\subsection{Técnicas de detección molecular}

\subsubsection{RT-qPCR}

La RT-qPCR combina transcripción reversa del ARN viral y amplificación cuantitativa de ADN complementario mediante PCR en tiempo real. En vigilancia de aguas residuales, esta técnica permite detectar material genético de virus de ARN y estimar la señal mediante valores de ciclo umbral o concentración calculada con curvas estándar. Su uso en SARS-CoV-2 fue central desde los primeros reportes de detección en aguas residuales y continúa siendo una herramienta de referencia para confirmar presencia viral antes de la secuenciación \cite{Ahmed_2020,Medema_2020,Parkins_2023}.

En matrices ambientales, la RT-qPCR requiere controles rigurosos. Los controles negativos permiten detectar contaminación; los controles positivos verifican desempeño del ensayo; los controles de extracción o recuperación evalúan pérdidas durante el procesamiento; y los marcadores fecales o de normalización ayudan a interpretar variaciones por dilución o carga de sólidos \cite{Parkins_2023,Hughes_2022}. En el estudio de RSV en sólidos sedimentados, Hughes y colaboradores incorporaron controles de recuperación, evaluación de inhibición y normalización con PMMoV, elementos relevantes para análisis de aguas residuales \cite{Hughes_2022}.

Las limitaciones principales de RT-qPCR en aguas residuales incluyen inhibidores de amplificación, degradación del ARN, baja concentración viral, heterogeneidad de la muestra y resultados cercanos al límite de detección. En la PTAR Quitumbe, algunas muestras con valores de Ct tardíos presentaron menor rendimiento para secuenciación, lo que muestra la relación práctica entre carga detectada por RT-qPCR y posibilidad de recuperar información genómica útil \cite{MejiaCalle_2024}. Por ello, la detección molecular debe interpretarse junto con controles de calidad y con las características del flujo de procesamiento de la muestra \cite{Child_2023,Parkins_2023}.

\subsubsection{PCR múltiple}

La PCR múltiple permite amplificar varios blancos genéticos dentro de una misma reacción o dentro de un flujo coordinado de amplificación. En vigilancia viral, puede emplearse para detectar distintos patógenos, diferenciar subtipos o generar conjuntos de amplicones que cubren regiones extensas del genoma. Su utilidad es alta cuando el volumen de muestra, el costo y el tiempo de procesamiento son variables críticas \cite{Child_2023,Le_2025}.

En secuenciación dirigida de virus, la amplificación múltiple mediante paneles de amplicones aumenta la probabilidad de recuperar información genómica a partir de muestras con ARN degradado o baja abundancia relativa. Child y colaboradores compararon métodos metagenómicos y dirigidos para virus humanos en aguas residuales y reportaron que la secuenciación por PCR de amplicones generó mayor cobertura para virus específicos que la metagenómica shotgun en muestras dominadas por material bacteriano \cite{Child_2023}. Para RSV, Le y colaboradores emplearon una PCR de amplicones solapados con 39 fragmentos de aproximadamente 550 pb, seguida de secuenciación Nanopore y análisis automático \cite{Le_2025}.

La interpretación de una PCR múltiple exige revisar blancos detectados, controles, curva o señal de amplificación, límites analíticos y coherencia entre réplicas. Resultados positivos con Ct tardíos requieren cautela cuando la carga es baja o la matriz contiene inhibidores; resultados negativos pueden reflejar ausencia del blanco, pérdida durante extracción, baja afinidad de cebadores o degradación del ácido nucleico \cite{Parkins_2023,Child_2023}. En análisis genómico, los sesgos de amplificación pueden favorecer regiones o linajes concretos y afectar la representación de variantes de baja frecuencia \cite{MejiaCalle_2024,Le_2025}.

\subsection{Secuenciación y caracterización genómica}

\subsubsection{Nanopore (ONT)}

La secuenciación Oxford Nanopore Technologies se basa en el paso de moléculas de ácido nucleico a través de nanopores y en la medición de cambios de corriente eléctrica para inferir bases en tiempo real. Esta tecnología permite lecturas largas, análisis portátil mediante dispositivos como MinION y generación rápida de datos, características útiles para vigilancia genómica aplicada a brotes, muestras ambientales y contextos con infraestructura variable \cite{Dash_2024,Shaw_2022}.

En vigilancia viral, ONT ha sido utilizado para SARS-CoV-2, poliovirus y RSV. En la PTAR Quitumbe, la secuenciación Nanopore permitió asignar linajes de SARS-CoV-2 desde muestras pasivas de aguas residuales y observar la presencia de variantes de Ómicron en el contexto local \cite{MejiaCalle_2024}. En poliovirus, Shaw y colaboradores describen el potencial de la detección molecular directa y la secuenciación Nanopore para reducir tiempos de confirmación de brotes \cite{Shaw_2022}. En RSV, Le y colaboradores emplearon ONT para genomas completos y análisis automático, con recuperación alta de cobertura en muestras con cargas superiores a 10.000 copias/mL \cite{Le_2025}.

La principal consideración técnica de ONT es su tasa de error cruda, generalmente mayor que la de plataformas de lectura corta. El control de calidad, la profundidad suficiente, el consenso, la depuración de lecturas y la selección adecuada de referencias reducen el impacto de estos errores en la interpretación final \cite{Dash_2024}. En aguas residuales, estos pasos adquieren mayor importancia por la presencia de mezclas virales, ARN fragmentado, material microbiano abundante y cobertura desigual entre amplicones \cite{Child_2023,MejiaCalle_2024}.

\subsubsection{Illumina}

La secuenciación Illumina se basa en síntesis por terminación reversible y produce lecturas cortas con alta precisión y gran profundidad. Es una de las plataformas más empleadas en secuenciación de nueva generación por su bajo error relativo, capacidad de multiplexación y utilidad para detectar variantes puntuales, reconstruir genomas de referencia y caracterizar diversidad genética en poblaciones virales \cite{Dash_2024}.

En aguas residuales, Illumina puede aplicarse mediante metagenómica shotgun, captura híbrida o paneles dirigidos. Child y colaboradores observaron que la metagenómica shotgun tuvo dificultad para recuperar cobertura suficiente de virus humanos en aguas residuales debido al predominio de ADN y ARN bacteriano; los métodos dirigidos y de enriquecimiento incrementaron la recuperación genómica de virus de interés \cite{Child_2023}. Esta evidencia indica que la plataforma de secuenciación debe elegirse junto con una estrategia de enriquecimiento compatible con la pregunta de investigación y con la carga viral esperada.

La fortaleza de Illumina reside en la precisión para identificar cambios de nucleótido y en la profundidad de lectura, mientras sus restricciones incluyen lecturas cortas, mayor dependencia de preparación de bibliotecas, requerimientos de equipo centralizado y necesidad de análisis bioinformático robusto. Para estudios de vigilancia ambiental, estas características la hacen valiosa como plataforma de confirmación, caracterización de diversidad y comparación con datos generados por tecnologías de lectura larga \cite{Dash_2024,Child_2023}.

\subsubsection{Variabilidad genética y vigilancia genómica}

La variabilidad genética viral permite estudiar transmisión, evolución, reemplazo de linajes, introducciones y señales mixtas en poblaciones. En SARS-CoV-2, la vigilancia genómica ha sido utilizada para identificar variantes regionalmente prevalentes en aguas residuales y para estimar mezclas de linajes en comunidades conectadas a sistemas de alcantarillado \cite{Crits_Christoph_2021,MejiaCalle_2024}. En RSV, la clasificación por subtipos y clados requiere referencias adecuadas para evitar asignaciones incompletas o pérdida de lecturas durante el mapeo \cite{Le_2025}. En poliovirus, la comparación genética de VP1 y genomas completos permite distinguir cepas Sabin-like, derivadas de vacuna y relaciones entre eventos de detección ambiental \cite{Zurbriggen_2008,Bottcher_2025}.

La vigilancia genómica complementa la detección molecular porque permite interpretar qué variantes o linajes sostienen una señal positiva. En aguas residuales, esta información tiene naturaleza poblacional: una misma muestra puede contener genomas procedentes de múltiples individuos, linajes y estados de infección. Por ello, la caracterización genómica aporta resolución epidemiológica cuando se acompaña de criterios de cobertura, profundidad, calidad de lectura, controles de contaminación y análisis de abundancia relativa \cite{Parkins_2023,Child_2023}.

La variabilidad detectada en una muestra ambiental también puede estar influida por procesos técnicos. La degradación del ARN, la eficiencia de concentración, la selección de cebadores, la competencia entre amplicones, la elección de referencia y el umbral de reporte pueden modificar la representación de linajes minoritarios \cite{MejiaCalle_2024,Le_2025}. Estos factores deben reconocerse al usar datos genómicos para vigilancia, especialmente cuando el objetivo es identificar cambios emergentes o señales de baja abundancia.

\subsection{Herramientas de análisis bioinformático}

\subsubsection{Procesamiento de secuencias}

El procesamiento bioinformático convierte lecturas crudas en información interpretable. En datos de secuenciación, las etapas generales incluyen demultiplexado, eliminación de adaptadores, filtrado por calidad, recorte, control de longitud, mapeo contra referencias, ensamblaje o generación de consenso, estimación de cobertura y revisión de profundidad por posición \cite{Dash_2024,Child_2023}. En ONT, el flujo incorpora basecalling y control de calidad de lecturas largas; en Illumina, se enfatiza la calidad por base, el recorte de adaptadores y el manejo de lecturas pareadas \cite{Dash_2024}.

Las muestras de aguas residuales agregan complejidad al procesamiento porque contienen mezclas de microorganismos, material humano, sólidos, inhibidores residuales y ácidos nucleicos degradados. En secuenciación viral, Child y colaboradores muestran que la elección entre metagenómica, captura híbrida y PCR dirigida cambia de forma importante la cobertura recuperada de virus humanos \cite{Child_2023}. En la PTAR Quitumbe, el uso de Freyja permitió estimar linajes de SARS-CoV-2 a partir de lecturas de aguas residuales, aunque el rendimiento dependió de la carga viral y de la calidad de amplificación \cite{MejiaCalle_2024}.

La trazabilidad del análisis es esencial para que los resultados sean reproducibles. Cada muestra debe conservar metadatos sobre fecha, punto de muestreo, método de concentración, controles, valores de Ct, plataforma de secuenciación, parámetros de filtrado y versión de herramientas utilizadas \cite{Parkins_2023,MejiaCalle_2024}. Sin estos elementos, la comparación temporal o entre patógenos pierde solidez metodológica.

\subsubsection{Identificación de variantes y linajes}

La identificación de variantes y linajes requiere comparar lecturas o consensos con bases de datos de referencia actualizadas. Para SARS-CoV-2, herramientas de asignación de linajes y estimación de abundancia permiten interpretar mezclas poblacionales en aguas residuales; Freyja fue utilizada en la PTAR Quitumbe para estimar linajes de Ómicron a partir de datos ambientales \cite{MejiaCalle_2024}. La experiencia de Crits-Christoph y colaboradores mostró que la secuenciación de aguas residuales puede identificar variantes regionalmente prevalentes, siempre que exista cobertura suficiente en posiciones informativas \cite{Crits_Christoph_2021}.

En RSV, la asignación requiere distinguir RSV-A y RSV-B y seleccionar referencias compatibles con el subtipo predominante. Le y colaboradores demostraron que el uso de una referencia inadecuada reduce el número de lecturas mapeadas, por lo que RSVTyper incorpora detección automática de subtipo y selección de referencia antes de análisis posteriores \cite{Le_2025}. En poliovirus, la identificación molecular se apoya en diferenciación intratípica, secuenciación de regiones como VP1 y comparación genética para reconocer cepas relacionadas o derivadas de vacuna \cite{Zurbriggen_2008,Asghar_2014,Bottcher_2025}.

Los criterios de confianza deben considerar profundidad mínima, cobertura del genoma o de regiones diagnósticas, concordancia entre réplicas, calidad de lectura, frecuencia relativa del linaje y consistencia con controles. En muestras ambientales con mezcla de genomas, la presencia de una variante minoritaria requiere umbrales explícitos para reducir reportes espurios por error de secuenciación, contaminación cruzada o mapeo ambiguo \cite{Dash_2024,Child_2023,MejiaCalle_2024}.

\subsubsection{Interpretación de resultados y limitaciones}

La interpretación de resultados en aguas residuales debe integrar variables analíticas, ambientales y epidemiológicas. Una señal positiva indica presencia de material genético compatible con el blanco analizado en la muestra, pero su magnitud puede variar por caudal, dilución, excreción individual, estabilidad del ARN, eficiencia de concentración y extracción, inhibición de PCR y cobertura poblacional del punto de muestreo \cite{Parkins_2023,Hughes_2022}. Por esta razón, las tendencias y comparaciones internas suelen ser más informativas que lecturas aisladas.

Las limitaciones genómicas incluyen baja concentración viral, ARN fragmentado, predominio de material bacteriano, sesgos de amplificación, cobertura incompleta, errores de lectura, mezcla de linajes y baja representación de variantes minoritarias. Child y colaboradores documentaron que la metagenómica shotgun puede ser insuficiente para recuperar genomas virales en aguas residuales sin enriquecimiento, mientras los métodos dirigidos aumentan cobertura a costa de depender de conocimiento previo del genoma \cite{Child_2023}. En el antecedente de Quitumbe, las muestras con Ct tardíos tuvieron menor éxito de secuenciación, y los cambios de cebadores afectaron la recuperación y asignación de linajes \cite{MejiaCalle_2024}.

La interpretación epidemiológica también debe reconocer que una planta de tratamiento representa una población conectada al sistema de alcantarillado, con variaciones por movilidad, cobertura de red, aportes no domésticos y cambios en el patrón de uso del agua. Los datos ambientales aportan evidencia de circulación y tendencias comunitarias, mientras la atribución a casos, barrios o grupos específicos requiere diseños adicionales y metadatos poblacionales robustos \cite{Parkins_2023,Sims_2020}.

\subsection{Contexto epidemiológico del estudio}

\subsubsection{Situación en Ecuador}

La evidencia local incorporada en este trabajo se concentra en el antecedente de vigilancia ambiental de SARS-CoV-2 realizado en la PTAR Quitumbe. Ese estudio documentó muestreo pasivo semanal de 24 horas, 47 muestras analizadas, cobertura de once barrios del Distrito Metropolitano de Quito y secuenciación de 36 muestras con Oxford Nanopore \cite{MejiaCalle_2024}. Sus resultados identificaron linajes de Ómicron y describieron la detección ambiental de JN.1 antes de su identificación clínica local, lo que sugiere valor del alcantarillado como sitio centinela para vigilancia genómica urbana \cite{MejiaCalle_2024}.

El contexto ecuatoriano descrito por Mejia Calle muestra una brecha de información genómica clínica, asociada con disponibilidad limitada de datos actualizados y dependencia de repositorios donde la carga de secuencias puede ser irregular \cite{MejiaCalle_2024}. Esta limitación coincide con observaciones regionales e internacionales sobre desafíos de vigilancia en contextos con recursos restringidos, donde la frecuencia de muestreo, el financiamiento, la infraestructura y la capacidad para monitorear variantes pueden condicionar la continuidad de los programas \cite{Parkins_2023}.

Para RSV y poliovirus, los recursos revisados ofrecen evidencia internacional sólida sobre vigilancia ambiental, vigilancia genómica y detección de circulación silenciosa, aunque aportan menos información específica de Ecuador. Esta situación refuerza la pertinencia de un estudio local que explore varios virus en una misma matriz ambiental de Quito. La PTAR Quitumbe constituye un punto de partida metodológicamente sustentado para ampliar la vigilancia desde SARS-CoV-2 hacia RSV y poliovirus, dos blancos con relevancia sanitaria y epidemiológica documentada \cite{Asghar_2014,Hughes_2022,Le_2025,Bottcher_2025}.

\subsubsection{Justificación del estudio}

El estudio se justifica por la necesidad de generar evidencia local sobre virus patógenos en aguas servidas de Quito, utilizando un sitio previamente caracterizado y una estrategia compatible con vigilancia ambiental. La PTAR Quitumbe reúne aportes de once barrios y ya demostró factibilidad para muestreo pasivo, detección molecular y secuenciación de SARS-CoV-2 \cite{MejiaCalle_2024}. Al incorporar RSV y poliovirus, el trabajo amplía el alcance hacia un enfoque multipatógeno respaldado por antecedentes internacionales de vigilancia en aguas residuales \cite{Asghar_2014,Hughes_2022,Bubba_2023}.

La importancia del estudio radica en que SARS-CoV-2, RSV y poliovirus representan problemas diferentes de vigilancia: variantes respiratorias de rápida sustitución, infecciones respiratorias con carga alta en grupos vulnerables y circulación entérica silenciosa relevante para erradicación. La vigilancia integrada puede generar información útil para investigación, salud pública y gestión de saneamiento, especialmente cuando se combinan detección molecular, secuenciación y análisis bioinformático \cite{Crits_Christoph_2021,Le_2025,Shaw_2022}.

La factibilidad se sostiene en tres elementos: un punto de muestreo local documentado, métodos moleculares aplicables a matrices ambientales y herramientas de secuenciación con antecedentes en virus humanos presentes en aguas residuales. Child y colaboradores muestran que los métodos dirigidos mejoran la recuperación genómica de virus específicos en matrices dominadas por material bacteriano; Le y colaboradores documentan un flujo ONT para RSV en muestras clínicas y de aguas residuales; y Shaw y colaboradores describen el valor potencial de Nanopore para acelerar la caracterización de poliovirus \cite{Child_2023,Le_2025,Shaw_2022}.

Los beneficiarios directos incluyen equipos académicos y técnicos que trabajan en microbiología ambiental, virología molecular, bioinformática, saneamiento y vigilancia epidemiológica. Los beneficiarios indirectos incluyen la población conectada a la PTAR Quitumbe, instituciones de salud pública y gestores urbanos que podrían utilizar evidencia ambiental para fortalecer sistemas de monitoreo en Quito. En conjunto, este marco teórico sostiene que el análisis molecular y genómico de SARS-CoV-2, RSV y poliovirus en aguas servidas puede aportar una base local para vigilancia viral ambiental, con alcances definidos y limitaciones metodológicas explícitas \cite{Parkins_2023,MejiaCalle_2024}.
